% Se explicita la indentación de 0 in para que el anexo
% no se vea perjudicado por el estilo de la bibliografía.
\leftskip 0in
\parindent 0.49cm % Sangría según el manual: la volvemos a establecer.
\newpage
\section{Anexos}

% =====================================================================
\subsection{Reflexión personal}
% =====================================================================
% [!! LÍMITE DURO: 2 PÁGINAS. En PRIMERA PERSONA. ~550-650 palabras.]
%
% [RÚBRICA C, 1 punto] Son solo TRES criterios y hay que cumplir los tres:
%   1. Reflexión sobre el aporte de la práctica al DESARROLLO PROFESIONAL.
%   2. Reflexión sobre la FORMACIÓN OTORGADA POR LA ESCUELA DE INGENIERÍA.
%   3. Estar escrita en PRIMERA PERSONA.
% El manual agrega un cuarto: destacar el grado de alcance de los objetivos de
% la Práctica II.
%
% >>> CAMBIO DE VOZ: este es el único lugar del informe donde escribes "yo".
%     Aprovéchalo. Después de 20 páginas en tercera persona, el contraste
%     funciona bien si la reflexión es concreta y no genérica.
%
% >>> REGLA ANTI-CLICHÉ: cada afirmación debe estar anclada a un hecho de TU
%     práctica. "Aprendí a trabajar en equipo" no vale nada. "Preparar las
%     reuniones con INVAP me obligó a anticipar las objeciones técnicas antes de
%     presentarlas" sí vale.
% =====================================================================

\subsubsection{Aporte al desarrollo profesional}
% [EXTENSIÓN] ~300 palabras, 3 párrafos.
%
% [MATERIAL CONCRETO DE TU PRÁCTICA -- elige 3 y desarróllalos]
%   - Fue mi primera experiencia asumiendo la responsabilidad completa de la
%     documentación de un proyecto: el CONOPS, la matriz de requerimientos y el
%     documento PDR llevan mi firma como elaborador.
%   - Las cuatro reuniones con INVAP: preparar preguntas técnicas de
%     caracterización de canal para una contraparte con décadas de experiencia
%     espacial me exigió llegar con el problema estudiado, no con dudas abiertas.
%   - La decisión de migrar el proyecto a ARCADIA/Capella la propuse yo, sin que
%     nadie me la pidiera. Aprender la metodología y migrar el proyecto completo
%     en una semana me mostró qué significa hacerse cargo de una decisión técnica.
%   - Descubrir el nulo cenital en la simulación: un resultado que contradecía mi
%     intuición inicial (más cerca no es mejor) y que me enseñó a desconfiar de
%     los modelos simplificados.
%   - Trabajar sin referencias internas: no había plantillas ni proyectos previos,
%     y eso cambió mi relación con la incertidumbre técnica.
%
% [ESTRUCTURA]
%   P1  Qué significó asumir la responsabilidad técnica y documental del proyecto
%       siendo practicante, y qué me costó al principio.
%   P2  Una situación concreta que me hizo crecer (elige UNA y cuéntala bien:
%       las reuniones con INVAP, o la migración a Capella).
%   P3  Qué me llevo para mi ejercicio profesional futuro.
%
% [EJEMPLO DE ARRANQUE DE P1]
%   "Cuando entré a DTS esperaba apoyar a un ingeniero con experiencia en un
%    proyecto ya definido. Lo que encontré fue un proyecto que no existía todavía
%    y un área sin desarrollos espaciales previos, de modo que la documentación
%    que produje no complementaba el trabajo de otro: era el trabajo. Firmar como
%    elaborador el documento de Revisión de Diseño Preliminar me obligó a
%    responder por cada decisión que había tomado en las ocho semanas anteriores."

% ESCRIBIR AQUÍ (primera persona)

\subsubsection{Formación otorgada por la Escuela de Ingeniería}
% [EXTENSIÓN] ~300 palabras, 3 párrafos.
%
% [QUÉ SÍ ME DIO LA ESCUELA -- conecta con cursos o competencias reales]
%   - Bases de comunicaciones y propagación: sin ellas no habría podido armar el
%     presupuesto de enlace ni entender el efecto del patrón de radiación.
%   - Procesamiento de señales y sistemas digitales: para entender la cadena
%     SDR, la demodulación en FPGA y la compensación Doppler.
%   - Programación y modelamiento en Python: la herramienta de simulación se
%     apoya directamente en eso.
%   - La costumbre de justificar cada decisión y de escribir informes técnicos.
%
% [QUÉ NO ME DIO Y TUVE QUE APRENDER SOLO -- esta parte es la que da valor a la
%  reflexión; sé honesto pero constructivo, no reclames]
%   - Ingeniería de sistemas como disciplina: requerimientos, trazabilidad,
%     revisiones de madurez, gestión de riesgos. No aparece en el plan de estudios
%     y fue el 80 % de mi práctica.
%   - Normativa y estándares técnicos: aprender a leer y aplicar un estándar
%     (IEEE, Argos, restricciones de plataforma) es una competencia distinta a
%     resolver un problema con enunciado cerrado.
%   - El entorno espacial: radiación, desgasificación, ciclado térmico.
%
% [OBLIGATORIO -- lo pide el manual] Cierra evaluando el grado de alcance de los
% objetivos de la Práctica II: aplicar conocimientos a un problema real de
% ingeniería, conocer la realidad laboral de la especialidad, comunicar de forma
% efectiva con equipos y superiores, y proponer soluciones a partir de
% herramientas modernas.
%
% [ESTRUCTURA]
%   P1  Qué de la formación de la Escuela usé directamente y en qué momento.
%   P2  Qué brecha encontré y cómo la cubrí durante la práctica.
%   P3  Balance del cumplimiento de los objetivos de la Práctica II.
%
% [EJEMPLO DE ARRANQUE DE P2]
%   "La brecha más grande que encontré no fue técnica sino metodológica. Durante
%    la carrera resolví problemas cuyo enunciado ya estaba escrito; en DTS mi
%    primera tarea fue escribir el enunciado. La ingeniería de sistemas
%    (la elicitación de requerimientos, la trazabilidad, las revisiones de
%    madurez) no forma parte del plan de estudios, y sin embargo ocupó la mayor
%    parte de mi trabajo."

% ESCRIBIR AQUÍ (primera persona)


% =====================================================================
%  ANEXOS TÉCNICOS
% =====================================================================
% Las tablas siguientes se generan automáticamente desde el Excel
% 07-TraceabilityRequirimentsMatrix.xlsx ejecutando:
%     python3 content/gen_tablas.py
% No las edites a mano: vuelve a correr el script si cambias el Excel.
% =====================================================================
\newpage
% ARCHIVO GENERADO AUTOMÁTICAMENTE por content/gen_tablas.py
% No editar a mano: vuelve a correr el script si cambia el Excel.

\subsection{Anexo A: requerimientos de usuario}

\begin{footnotesize}
\begin{longtable}{@{}l >{\raggedright\arraybackslash}p{5.4cm} c c c >{\raggedright\arraybackslash}p{4.6cm}@{}}
\caption{Requerimientos de usuario del proyecto BIFROST, con su prioridad, nivel de riesgo, tipo y mitigación definida. Elaboración propia a partir de la matriz de trazabilidad del proyecto.\vspace{1ex}}\label{tab:anexoA}\\
\toprule
\textbf{ID} & \textbf{Requerimiento} & \textbf{Prior.} & \textbf{Riesgo} & \textbf{Tipo} & \textbf{Mitigación propuesta} \\
\midrule\endfirsthead
\multicolumn{6}{@{}l@{}}{\textit{(continuación)}}\\
\toprule
\textbf{ID} & \textbf{Requerimiento} & \textbf{Prior.} & \textbf{Riesgo} & \textbf{Tipo} & \textbf{Mitigación propuesta} \\
\midrule\endhead
\bottomrule\endfoot
UR-01 & El sistema deberá proveer un servicio Store \& Forward satelital, almacenando datos durante su orbita para ser descargados en una estación terrestre. & Alta & Alto & Funcional & Diseño de memoria no volátil con margen $\geq$50\%, simulación de carga orbital y pruebas de estrés en tierra \\
\addlinespace[2pt]
UR-02 & El sistema deberá soportar un enlace downlink desde el satélite hacia estaciones terrenas en territorio chileno. & Alta & Medio & Funcional &  \\
\addlinespace[2pt]
UR-03 & El sistema deberá soportar un enlace uplink desde estaciones remotas de monitoreo meteorológico, sismológico y/o vulcanológico hacia el satélite para la recolección de datos. & Alta & Medio & Funcional & Validación temprana con transmisores representativos y pruebas HIL \\
\addlinespace[2pt]
UR-04 & El sistema deberá ser compatible con el estándar internacional para recolección de datos ambientales por satélite utilizado por estaciones remotas autónomas de meteorología, sismología, vulcanología, etc. & Alta & Medio & Proceso & Revisión de estándar ARGOS y asesoría INVAP \\
\addlinespace[2pt]
UR-05 & El sistema deberá asegurar la integridad de los datos durante su recepción, almacenamiento y transmisión. & Media & Bajo & Funcional &  \\
\addlinespace[2pt]
UR-06 & El sistema deberá ser compatible con plataforma CubeSat 12U en órbita terrestre baja (LEO) sol síncrona a una altitud nominal aproximada de 650 km con una inclinación de 96$^{\circ}$. & Alta & Alto & No funcional & Análisis de interfaces mecánicas/eléctricas y revisión PDR con proveedor de plataforma \\
\addlinespace[2pt]
UR-07 & El sistema deberá cubrir el territorio chileno continental, zonas extremas, territorio insular y territorio antártico. & Media & Bajo & Funcional &  \\
\addlinespace[2pt]
UR-08 & El sistema deberá permitir el uso de componentes COTS siempre que cumplan con: (a) funcionamiento en órbita probada, (b) tolerancia a la radiación verificada, o (c) mitigaciones de diseño que reduzcan el riesgo a un nivel aceptable para la misión. & Alta & Alto & No funcional & Selección COTS con validación espacial y mitigación por redundancia y watchdogs \\
\addlinespace[2pt]
UR-09 & El sistema deberá ser de carácter experimental, con proyección a futuras aplicaciones comerciales. & Baja & Bajo & No funcional &  \\
\addlinespace[2pt]
UR-10 & El sistema deberá presentar una arquitectura modular, reutilizable y escalable. & Baja & Bajo & No funcional &  \\
\addlinespace[2pt]
UR-11 & El sistema debe garantizar una vida útil operativa mínima de 5 años en el entorno de radiación LEO definido. & Alta & Alto & No funcional & Pruebas de estrés, análisis térmico y pruebas de vida aceleradas \\
\addlinespace[2pt]
UR-12 & El sistema debe ser validado mediante pruebas funcionales en tierra y demostración en órbita. & Baja & Bajo & Proceso &  \\
\addlinespace[2pt]
UR-13 & El sistema deberá ser compatible con el sistema de separación del lanzador seleccionado EXOpod 12U, cumpliendo con las restricciones geométricas, mecánicas y de seguridad. &  &  & No funcional &  \\
\addlinespace[2pt]
\end{longtable}
\end{footnotesize}

\subsection{Anexo B: requerimientos de sistema}

\begin{footnotesize}
\begin{longtable}{@{}l >{\raggedright\arraybackslash}p{5.4cm} c c c >{\raggedright\arraybackslash}p{4.6cm}@{}}
\caption{Requerimientos de sistema derivados de los requerimientos de usuario, con su prioridad, nivel de riesgo, tipo y mitigación definida. Elaboración propia a partir de la matriz de trazabilidad del proyecto.\vspace{1ex}}\label{tab:anexoB}\\
\toprule
\textbf{ID} & \textbf{Requerimiento} & \textbf{Prior.} & \textbf{Riesgo} & \textbf{Tipo} & \textbf{Mitigación propuesta} \\
\midrule\endfirsthead
\multicolumn{6}{@{}l@{}}{\textit{(continuación)}}\\
\toprule
\textbf{ID} & \textbf{Requerimiento} & \textbf{Prior.} & \textbf{Riesgo} & \textbf{Tipo} & \textbf{Mitigación propuesta} \\
\midrule\endhead
\bottomrule\endfoot
SR-01 & El sistema deberá permitir la recepción de paquetes de datos con una tamaño desde 32 a 256 bits, con una tasa de datos de 400 bps. & Alta & Alto & Funcional & Prototipo temprano del receptor y pruebas BER en laboratorio \\
\addlinespace[2pt]
SR-02 & El sistema deberá operar como receptor en el rango de frecuencias 401 - 403 MHz utilizado por estaciones remotas, con una polarización circular derecha y un ancho de haz de 130$^{\circ}$ HPBW. & Alta & Alto & No funcional & Verificación de front-end RF y filtros certificados \\
\addlinespace[2pt]
SR-03 & El sistema deberá ser compatible con la estructura de datos enviado por las estaciones remotas , incluyendo modulación, canalización y protocolos, además de las variaciones en frecuencia por efecto Doppler. & Alta & Alto & Funcional & Simulación Doppler y pruebas con generador de señales \\
\addlinespace[2pt]
SR-04 & El sistema deberá almacenar al menos 1.000.000 mensajes de durante un tiempo máximo de 15 órbitas y sin pérdidas hasta establecer contacto con una estación terrestre para la descarga. & Media & Medio & Funcional & Margen de memoria $\geq$50\% y control de overflow \\
\addlinespace[2pt]
SR-05 & El sistema deberá mantener la asociación temporal de los datos almacenados. & Baja & Bajo & Funcional &  \\
\addlinespace[2pt]
SR-06 & El sistema deberá permitir la descarga de al menos 10.000 mensajes almacenados hacia estaciones terrestres cuando exista conectividad disponible. & Media & Medio & Funcional & Ensayos de descarga en peor escenario \\
\addlinespace[2pt]
SR-07 & El sistema deberá entregar los datos al subsistema de comunicaciones de la plataforma (Transponder Banda S) para el enlace downlink. & Media & Medio & Proceso &  \\
\addlinespace[2pt]
SR-08 & El sistema deberá comunicarse con el OBC del CubeSat 12U mediante el bus definido por la plataforma (CANBUS), gestionando una tasa de transferencia de datos compatible y reportando periódicamente su estado operativo. & Alta & Alto & No funcional &  \\
\addlinespace[2pt]
SR-09 & El sistema deberá operar bajo el control del OBC, escuchando activamente en las zonas definidas por la misión y transmitiendo únicamente cuando sea ordenado. & Media & Bajo & Funcional &  \\
\addlinespace[2pt]
SR-10 & El sistema deberá implementar mecanismos de integridad de datos en todas las etapas, con un error de bit menor a 0.01\% y control de acceso a la descarga de datos por estaciones terrestres. & Media & Medio & Interfaz & CRC, ACK/NACK y control de acceso \\
\addlinespace[2pt]
SR-13 & El sistema deberá consumir un máximo de 10 W en operación nominal y ser alimentado por el sistema de potencia del satélite de forma continua. & Alta & Alto & No funcional & Pruebas de compatibilidad CAN con OBC real \\
\addlinespace[2pt]
SR-14 & El sistema deberá tener dimensiones máximas 20 cm x 20 cm x 10 cm. & Alta & Alto & No funcional & Validación de lógica de estados y temporización \\
\addlinespace[2pt]
SR-15 & El sistema deberá tener una masa máxima de 3.0 kg y mantener su Centro de Gravedad (CoG) a $\pm$2 mm del centro geométrico, el cual coincide con el eje logitudinal geométrico del satélite. & Alta & Alto & No funcional & Presupuesto de potencia con margen $\geq$25\% \\
\addlinespace[2pt]
SR-16 & El sistema deberá cumplir con requisitos de EMI/EMC compatibles con operación satelital, evitando interferencias internas y externas durante todo el proceso de lanzamiento, despliege y operación. & Alta & Alto & No funcional & Revisión de layout 3D y mockup mecánico \\
\addlinespace[2pt]
SR-17 & El sistema deberá diseñarse de forma modular y permitir la actualización de firmware en vuelo (OTA). & Media & Bajo & No funcional & Simulaciones CAD para estimaciones de masas y revisión periodica de masa en el prototipo \\
\addlinespace[2pt]
SR-18 & El sistema deberá certificarse mediante prueba térmo-vacio, mecánicas, EMI y eléctricas en las instalaciones de Centro Espacial Nacional y ser validado en órbita. & Media & Bajo & Proceso & Mediciones EMI/EMC y separación RF \\
\addlinespace[2pt]
SR-19 & El sistema deberá ser respetar la envolvente interna permita, los volúmen prohibídos y las restricciones sobre protrusiones, mecanismos móviles y bordes activos, asegurando la integración con la plataforma. &  &  & No funcional &  \\
\addlinespace[2pt]
SR-20 & El sistema deberá soportar los ambientes mecánicos de lanzamiento (cargas cuasiestáticas, vibración y shock) definidos por el lanzador y deployer seleccionados. &  &  & No funcional &  \\
\addlinespace[2pt]
SR-21 & El sistema deberá implementar Watchdog Timers externos e independientes para reinicio ante eventos SEU/Latch-up. &  &  & Funcional &  \\
\addlinespace[2pt]
SR-22 & Protección electrónica contra Single Event Latch-up (SEL) mediante limitación de corriente activa. &  &  & No funcional &  \\
\addlinespace[2pt]
SR-23 & Los componentes y materiales no metálicos que uitlice el sistema deberán cumplir con un TML $<$ 1.0\% y un CVCM $<$0.1\% &  &  & No funcional &  \\
\addlinespace[2pt]
\end{longtable}
\end{footnotesize}

\subsection{Anexo C: matriz de trazabilidad de requerimientos}

\begin{footnotesize}
\begin{longtable}{@{}l >{\raggedright\arraybackslash}p{12.2cm}@{}}
\caption{Trazabilidad entre requerimientos de usuario y requerimientos de sistema. Elaboración propia.\vspace{1ex}}\label{tab:anexoC}\\
\toprule
\textbf{Req. de usuario} & \textbf{Requerimientos de sistema asociados} \\
\midrule\endfirsthead
\multicolumn{2}{@{}l@{}}{\textit{(continuación)}}\\
\toprule
\textbf{Req. de usuario} & \textbf{Requerimientos de sistema asociados} \\
\midrule\endhead
\bottomrule\endfoot
UR-01 & SR-04, SR-06 \\
\addlinespace[2pt]
UR-02 & SR-03, SR-07 \\
\addlinespace[2pt]
UR-03 & SR-02, SR-03 \\
\addlinespace[2pt]
UR-04 & SR-01, SR-02, SR-03 \\
\addlinespace[2pt]
UR-05 & SR-04, SR-05, SR-08, SR-10 \\
\addlinespace[2pt]
UR-06 & SR-13, SR-14, SR-15, SR-16, SR-19, SR-20 \\
\addlinespace[2pt]
UR-07 & SR-09, SR-10 \\
\addlinespace[2pt]
UR-08 & SR-16, SR-17, SR-18, SR-21, SR-23 \\
\addlinespace[2pt]
UR-09 & SR-17 \\
\addlinespace[2pt]
UR-10 & SR-17 \\
\addlinespace[2pt]
UR-11 & SR-17, SR-18, SR-21, SR-23 \\
\addlinespace[2pt]
UR-12 & SR-18 \\
\addlinespace[2pt]
UR-13 & SR-19, SR-22 \\
\addlinespace[2pt]
\end{longtable}
\end{footnotesize}

\subsection{Anexo D: matriz preliminar de ensayos de verificación}

\begin{footnotesize}
\begin{longtable}{@{}l >{\raggedright\arraybackslash}p{4.0cm} >{\raggedright\arraybackslash}p{8.6cm}@{}}
\caption{Ensayos de verificación especificados para las fases de diseño detallado y calificación del sistema. Elaboración propia.\vspace{1ex}}\label{tab:anexoD}\\
\toprule
\textbf{ID} & \textbf{Nombre del ensayo} & \textbf{Descripción} \\
\midrule\endfirsthead
\multicolumn{3}{@{}l@{}}{\textit{(continuación)}}\\
\toprule
\textbf{ID} & \textbf{Nombre del ensayo} & \textbf{Descripción} \\
\midrule\endhead
\bottomrule\endfoot
FUN-01 & Caracterización de Potencia & Medición de consumo de corriente en todos los modos (Sleep, RX, Processing, TX Interface). Verificación de picos de arranque (Inrush Current). \\
\addlinespace[2pt]
FUN-02 & Sensibilidad y Selectividad RF & Verificación del umbral de recepción (-115 dBm estimado) y ancho de banda en 401-403 MHz usando generador de señales. \\
\addlinespace[2pt]
FUN-03 & Simulación de Efecto Doppler & Inyección de señales con desplazamiento de frecuencia variable ($\pm$10 kHz) para validar el algoritmo de corrección del receptor. \\
\addlinespace[2pt]
FUN-04 & Integridad de Datos (BER) & Transmisión de 1 millón de paquetes y verificación de tasa de error de bit (BER) y corrección CRC tras almacenamiento. \\
\addlinespace[2pt]
FUN-05 & Interfaz CAN Bus \& CSP & Validación de intercambio de comandos y telemetría con emulador de OBC (o OBC real). Verificación de manejo de errores de bus. \\
\addlinespace[2pt]
FUN-06 & Prueba de Watchdog y Failsafe & Inyección de fallos de software y hardware (reset forzado) para medir tiempo de recuperación y estado seguro. \\
\addlinespace[2pt]
FUN-07 & Actualización de Firmware OTA & Simulación de carga de nueva imagen de software a través del bus de datos y verificación de checksum e instalación. \\
\addlinespace[2pt]
THR-01 & Ciclado Térmico (Presión Ambiente) & Ciclos de temperatura entre -30$^{\circ}$C y +70$^{\circ}$C en cámara climática para detectar fallas de soldadura o componentes infantiles (Stress screening). \\
\addlinespace[2pt]
THR-02 & Vacío Térmico (TVAC) - Bake-out & Exposición a alto vacío ($<$10$^{-5}$ Torr) y temperatura elevada (+60$^{\circ}$C) para desgasificación (Outgassing) y limpieza. \\
\addlinespace[2pt]
THR-03 & Balance Térmico en Vacío & Validación del modelo térmico: verificación de temperaturas de operación en vacío (solo conducción) en casos fríos y calientes. \\
\addlinespace[2pt]
MEC-01 & Propiedades de Masa & Medición precisa de masa total ($<$ 3kg) y determinación del Centro de Gravedad (CoG) y Momentos de Inercia (MoI). \\
\addlinespace[2pt]
MEC-02 & Vibración Aleatoria (Random Vib) & Ensayo en shaker electrodinámico según perfiles de lanzamiento (ej. NASA GEVS o perfil Falcon 9). Verificación de integridad estructural. \\
\addlinespace[2pt]
MEC-03 & Vibración Sinusoidal (Sine Sweep) & Barrido de frecuencias para identificar resonancias naturales (Eigenfrequencies) y asegurar que $>$ 100 Hz (típico CubeSat). \\
\addlinespace[2pt]
MEC-04 & Ensayo de Choque (Shock) & Simulación de cargas de choque por separación (si aplica por requerimiento de lanzador) o despliegue de antenas. \\
\addlinespace[2pt]
MEC-05 & Despliegue de Antenas & Verificación funcional del mecanismo de despliegue de antenas UHF (Burn-wire / Resorte) en condiciones de gravedad compensada. \\
\addlinespace[2pt]
EMC-01 & Emisiones Conducidas (CE) & Medición de ruido inyectado por la carga útil hacia las líneas de potencia del bus (Power Bus Noise). \\
\addlinespace[2pt]
EMC-02 & Susceptibilidad Radiada (RS) & Verificación de operación nominal ante campos externos (ej. Radares de lanzamiento o transmisor Banda S propio). \\
\addlinespace[2pt]
EMC-03 & Emisiones Radiadas (RE) & Medición de ruido RF generado por la electrónica digital (relojes, CPU) para asegurar "silencio" en banda GPS y Banda S receptor. \\
\addlinespace[2pt]
EMC-04 & Auto-Compatibilidad & Verificación de que la transmisión UHF propia no resetea la CPU ni afecta sensores cercanos. \\
\addlinespace[2pt]
\end{longtable}
\end{footnotesize}

